% =====================================================================
%  Bölüm 4 — Yapıcı ve Klasik Mantık
% =====================================================================
\chapter{Yapıcı ve Klasik Mantık}\label{bol:yapici-klasik}

\begin{bolumhedefleri}
\begin{itemize}
  \item Brouwer--Heyting--Kolmogorov (BHK) yorumunu ve ``kanıt bir
        inşadır'' fikrini,
  \item \lc{Classical.em} ve \lc{Classical.byContradiction} ile klasik
        ispat yapmayı,
  \item Üçüncü hâlin olanaksızlığı, çift değilleme eleme ve Peirce
        yasasının birbirine denk olduğunu,
  \item \lc{\#print axioms} ile bir ispatın klasik olup olmadığını
        denetlemeyi,
  \item Bir Kripke modeliyle $P \vee \neg P$'nin neden yapıcı olarak
        ispatlanamadığını,
  \item Glivenko teoremini, çift değilleme çevirisini ve karar
        verilebilir önermeleri.
\end{itemize}
\end{bolumhedefleri}

Önceki iki bölümde \Lean{}'in bazı önermeleri ispatlamamıza izin
vermediğini gördük: Peirce yasası (Alıştırma~\ref{alis:peirce}),
De Morgan'ın bir yönü, $\neg\neg P \to P$ ve $P \vee \neg P$. Bu bölüm
bu sınırın ne olduğunu, neden var olduğunu ve nasıl aşılacağını
anlatıyor.

% ---------------------------------------------------------------------
\section{İki kanıt anlayışı}\label{sec:iki-anlayis}

\subsection{Yapıcı olmayan bir ispat}

Şu teoremi ve onun ünlü ispatını ele alalım.

\begin{teorem}{}{irrasyonel-us}
$a^b$ rasyonel olacak şekilde $a$ ve $b$ irrasyonel sayıları vardır.
\end{teorem}

\noindent\textit{İspat.} $x = \sqrt{2}^{\sqrt{2}}$ sayısına bakalım.
Bu sayı ya rasyoneldir ya değildir.
\begin{itemize}
  \item $x$ rasyonelse, $a = b = \sqrt{2}$ alırız. İşimiz biter.
  \item $x$ rasyonel değilse, $a = x$ ve $b = \sqrt{2}$ alırız. O zaman
        $a^b = \bigl(\sqrt{2}^{\sqrt{2}}\bigr)^{\sqrt{2}} = \sqrt{2}^{\,2} = 2$
        olur ve bu rasyoneldir. \hfill\qedsymbol
\end{itemize}

İspat kısa, zarif ve kusursuzdur, ama tuhaf bir eksiği vardır: Sonunda
$a$'nın ve $b$'nin \emph{ne olduğunu} hâlâ bilmiyoruz. İspat, ``ya
$(\sqrt2, \sqrt2)$ ya da $(\sqrt2^{\sqrt2}, \sqrt2)$ çiftinin işe yaradığını''
gösterir, ama hangisinin işe yaradığını söylemez. Bunu söyleyemez, çünkü
ilk adımdaki ``ya rasyoneldir ya değildir'' cümlesi, üçüncü hâlin
olanaksızlığının bir örneğidir ve hangisinin doğru olduğu hakkında hiçbir
bilgi taşımaz.

\begin{tarihce}[Tarihçe: Brouwer ve Hilbert]
Hollandalı matematikçi L.\,E.\,J. Brouwer, 1908'de yayımladığı bir
makalede mantık ilkelerinin, özellikle de üçüncü hâlin olanaksızlığının,
sonsuz kümeler üzerinde güvenilir olmadığını ileri sürdü
\citep{brouwer1908}. Brouwer'e göre matematiksel bir nesnenin var
olduğunu söylemek, onu zihinde \emph{inşa edebilmek} demekti. Bu görüş
\emph{sezgicilik} (İng.\ \emph{intuitionism}) adını aldı.

David Hilbert bu programa sert biçimde karşı çıktı. Ona göre üçüncü hâlin
olanaksızlığını matematikçinin elinden almak, gökbilimcinin elinden
teleskobunu ya da boksörün elinden yumruklarını almak gibiydi. 1920'lerde
yaşanan bu tartışma, matematiğin temelleri üzerine yirminci yüzyılın en
ünlü tartışmalarından biri oldu. Bugün ise iki taraf barışık
yaşıyor: Yapıcı mantık, hesaplanabilirliği gözeten bir \emph{alt sistem}
olarak, klasik mantık ise onun bir \emph{uzantısı} olarak görülüyor.
\end{tarihce}

\subsection{BHK yorumu}

Brouwer'in fikri, Arend Heyting ve Andrey Kolmogorov tarafından kesin bir
biçime kavuşturuldu \citep{heyting1930, kolmogorov1932}. Bugün
\emph{Brouwer--Heyting--Kolmogorov (BHK) yorumu} olarak bilinen bu
yaklaşım, her bağlaç için bir kanıtın \emph{ne olduğunu} söyler
(\cref{tab:bhk}).

\begin{table}[htbp]
\centering
\renewcommand{\arraystretch}{1.5}
\rowcolors{2}{tablosatir}{white}
\begin{tabular}{>{\centering\arraybackslash}p{1.7cm}
                >{\raggedright\arraybackslash}p{6.2cm}
                >{\raggedright\arraybackslash}p{3.6cm}}
\rowcolor{tabloust}
\tablobaslik{Önerme} & \tablobaslik{Bir kanıtı şudur…} & \tablobaslik{Lean'de} \\
$P \wedge Q$      & $P$'nin bir kanıtı \emph{ve} $Q$'nun bir kanıtından oluşan bir çift & \lc{⟨hp, hq⟩} \\
$P \vee Q$        & $P$'nin bir kanıtı \emph{ya da} $Q$'nun bir kanıtı, ve \textbf{hangisi olduğunu gösteren bir etiket} & \lc{Or.inl hp}, \lc{Or.inr hq} \\
$P \to Q$         & $P$'nin her kanıtını $Q$'nun bir kanıtına dönüştüren bir \emph{yöntem} & \lc{fun hp => ...} \\
$\neg P$          & $P$'nin her kanıtını bir çelişkiye dönüştüren bir yöntem & \lc{fun hp => ...} \\
$\bot$            & kanıtı yoktur & --- \\
\end{tabular}
\caption{BHK yorumu. Önceki iki bölümde Lean'de yaptığımız her şey bu
tabloya uyar.}
\label{tab:bhk}
\end{table}

Tablonun kalın yazılmış satırı her şeyi açıklar. BHK yorumunda
$P \vee \neg P$'nin bir kanıtı, ya $P$'nin ya $\neg P$'nin bir kanıtı
olmalıdır, \emph{ve hangisi olduğu belli olmalıdır}. Her önerme için
böyle bir kanıt vermek, matematikteki her sorunu çözebilen genel bir
yöntem bulmak demektir. Böyle bir yöntem yoktur (bu olguyu
\ciltiki{6}'da Turing'in sonuçlarıyla kesinleştireceğiz). Bu
yüzden $P \vee \neg P$, BHK yorumunda genel bir yasa olamaz.

\begin{onemlifikir}[Önemli Fikir: Yapıcı mantık daha az şey söylemez, daha çok şey söyler]
Yapıcı bir ispat, klasik bir ispattan \emph{daha fazla bilgi} taşır.
``$a^b$ rasyonel olan irrasyonel $a, b$ vardır'' cümlesinin yapıcı bir
ispatı, bu sayıları bize \emph{verir}. Klasik ispat ise yalnızca
var olduklarını söyler. Bu nedenle yapıcı mantıkta ispatlanan her şey
klasik mantıkta da doğrudur, ama tersi geçerli değildir.
\end{onemlifikir}

% ---------------------------------------------------------------------
\section{Lean'de klasik mantık}\label{sec:lean-klasik}

\Lean{}'de klasik mantık yasaklanmamıştır; yalnızca \emph{açıkça
istenmesi} gerekir. Kapı, \lc{Classical} ad alanındaki iki teoremdir:

\leandosya[Klasik mantığın iki kapısı]{Bolum04}{em_check}

\subsection{Çift değilleme eleme}

\lc{Classical.em P}, \lc{P ∨ ¬P} türünde bir terimdir. Ona \leantac{cases}
uyguladığımızda, iki durum açılır: $P$ doğru ya da $\neg P$ doğru. Bu
\emph{klasik durum ayırma}, yapıcı mantıkta mümkün olmayan şeylerin
kapısını açar:

\leandosya[Çift değilleme eleme]{Bolum04}{cift_degil_eleme}

\begin{figure}[htbp]
\centering
\begin{tikzpicture}[
    durum/.style={draw=#1, fill=#1!8, thick, rounded corners=4pt,
                  minimum width=3.4cm, minimum height=0.9cm, align=center,
                  font=\small},
    ok/.style={-{Stealth[length=2.6mm]}, thick, #1},
  ]
  \node[durum=teoremframe] (em) at (0,0) {\lc{Classical.em P : P ∨ ¬P}};
  \node[durum=ikincil] (a) at (-3.3,-1.7) {\lc{hp : P}};
  \node[durum=vurgu]   (b) at (3.3,-1.7)  {\lc{hnp : ¬P}};
  \node[durum=ikincil] (a2) at (-3.3,-3.3) {\lc{exact hp}};
  \node[durum=vurgu]   (b2) at (3.3,-3.3)  {\lc{hnnp hnp : False}\\[-1pt]{\scriptsize\lc{absurd hnp hnnp}}};
  \node[durum=bittiframe] (r) at (0,-4.9) {\lc{P} kanıtlandı};
  \draw[ok=ikincil] (em) -- node[left=6pt, font=\leanfont\scriptsize\bfseries, text=ikincil]{inl} (a);
  \draw[ok=vurgu]   (em) -- node[right=6pt, font=\leanfont\scriptsize\bfseries, text=vurgu]{inr} (b);
  \draw[ok=ikincil] (a) -- (a2);  \draw[ok=vurgu] (b) -- (b2);
  \draw[ok=bittiframe] (a2) -- (r); \draw[ok=bittiframe] (b2) -- (r);
  \node[font=\baslikfont\scriptsize, text=teoremframe, align=center] at (0,0.95)
       {klasik mantık burada devreye girer};
\end{tikzpicture}
\caption{\lc{Classical.em} ile durum ayırma. İkinci durumda elimizdeki
\lc{¬¬P} ile \lc{¬P} çelişir.}
\label{fig:klasik-durum}
\end{figure}

\subsection{Çelişki ile ispat}

Klasik matematiğin en sık kullanılan yöntemlerinden biri \emph{çelişki
ile ispattır} (Latince \emph{reductio ad absurdum}): ``$Q$'yu göstermek
için $Q$'nun yanlış olduğunu varsayalım ve bir çelişkiye ulaşalım.''

\leandosya[Çelişki ile ispat]{Bolum04}{celiski_ile_ispat}

\begin{dikkat}[Dikkat: İki tür ``çelişki ile ispat'']
Matematikçiler ``çelişki ile ispat'' derken iki farklı şeyi kastedebilir:
\begin{enumerate}[label=(\alph*)]
  \item \textbf{$\neg P$'yi göstermek için} $P$'yi varsayıp çelişkiye
        ulaşmak. Bu, değillemenin \emph{tanımıdır} ($\neg$G kuralı) ve
        tamamen yapıcıdır. $\sqrt{2}$'nin rasyonel \emph{olmadığının}
        ispatı bu türdendir.
  \item \textbf{$P$'yi göstermek için} $\neg P$'yi varsayıp çelişkiye
        ulaşmak. Bu, aslında $\neg\neg P$'yi gösterip ardından çift
        değillemeyi elemektir ve \emph{klasik} bir adımdır.
\end{enumerate}
\Lean{}'de (a) \leantac{intro} ile, (b) ise
\lc{Classical.byContradiction} ile yapılır. İkisini ayırt etmek, bir
ispatın gerçekten klasik mantığa ihtiyaç duyup duymadığını anlamanın
ilk adımıdır.
\end{dikkat}

\subsection{Eksik kalan ispatlar}

Önceki bölümlerde takıldığımız iki önermeyi artık ispatlayabiliriz:

\leandosya[De Morgan, zor yön]{Bolum04}{demorgan_ve_klasik}
\leandosya[Peirce yasası]{Bolum04}{peirce}

\noindent İkisinde de yöntem aynıdır: \lc{Classical.em P} ile durum ayır,
iki durumda da hedefe ulaş. Bu iki ispat, Bölüm~2'de açık bıraktığımız
Peirce yasası sorusunun cevabıdır: Peirce yasası \emph{klasik} bir
yasadır.

% ---------------------------------------------------------------------
\section{Üç ilke, tek sınır}\label{sec:uc-ilke}

Şimdiye kadar klasik mantığa özgü üç ilkeyle karşılaştık. Bunları
\emph{bütün önermeler için} söylenen birer önerme olarak
tanımlayalım:

\leandosya[Üç klasik ilke]{Bolum04}{ilkeler}

\begin{teorem}{Klasik ilkelerin denkliği}{uc-ilke}
Yapıcı mantıkta aşağıdaki üç ilke birbirine denktir:
\begin{center}
\begin{tabular}{rl}
\textbf{EM} (üçüncü hâlin olanaksızlığı): & $\forall P,\; P \vee \neg P$ \\
\textbf{DNE} (çift değilleme eleme):       & $\forall P,\; \neg\neg P \to P$ \\
\textbf{Peirce yasası}:                    & $\forall P\, Q,\; ((P \to Q) \to P) \to P$
\end{tabular}
\end{center}
Yani yapıcı mantığa bunlardan herhangi birini eklemek, ötekileri de
eklemek demektir. Ortaya çıkan sistem klasik mantıktır.
\end{teorem}

\noindent İspatı bir döngü olarak kuracağız (\cref{fig:uc-ilke}):
EM $\Rightarrow$ DNE $\Rightarrow$ Peirce $\Rightarrow$ EM.

\begin{figure}[htbp]
\centering
\begin{tikzpicture}[
    ilke/.style={draw=teoremframe, fill=teorembg, very thick, rounded corners=5pt,
                 minimum width=2.6cm, minimum height=1.1cm, align=center,
                 font=\baslikfont\bfseries},
    ok/.style={-{Stealth[length=3mm]}, very thick, anarenk},
    etiket/.style={font=\leanfont\scriptsize, text=anarenk, fill=white, inner sep=2pt},
  ]
  \node[ilke] (em)  at (90:2.3)  {EM\\[-1pt]\normalfont\scriptsize $P \vee \neg P$};
  \node[ilke] (dne) at (-30:2.3) {DNE\\[-1pt]\normalfont\scriptsize $\neg\neg P \to P$};
  \node[ilke] (pe)  at (210:2.3) {Peirce\\[-1pt]\normalfont\scriptsize $((P{\to}Q){\to}P){\to}P$};
  \draw[ok] (em) to[bend left=22] node[etiket, right=2pt]{em\_dne} (dne);
  \draw[ok] (dne) to[bend left=22] node[etiket, below=2pt]{dne\_peirce} (pe);
  \draw[ok] (pe) to[bend left=22] node[etiket, left=2pt]{peirce\_em} (em);
  \node[font=\baslikfont\small\bfseries, text=koyu, align=center] at (0,-0.1)
       {hepsi\\yapıcı};
\end{tikzpicture}
\caption{Üç klasik ilkenin denkliği. Her ok, \emph{klasik mantık
kullanmadan} ispatlanmış bir \Lean{} teoremidir.}
\label{fig:uc-ilke}
\end{figure}

\leandosya[EM $\Rightarrow$ DNE]{Bolum04}{em_dne}
\leandosya[DNE $\Rightarrow$ Peirce]{Bolum04}{dne_peirce}
\leandosya[Peirce $\Rightarrow$ EM]{Bolum04}{peirce_em}

\noindent Son ispattaki hileye dikkat edin: Peirce yasasındaki $Q$
herhangi bir önerme olabileceği için \lc{False}'u seçtik. O zaman
\lc{P ∨ ¬P → False}, yani \lc{¬(P ∨ ¬P)} elimize geçer ve
\cref{sec:em-yanlis-olamaz}'deki ``fikir değiştirme'' hilesi işi
bitirir. Döngü tamamlanınca denklikler hemen çıkar:

\leandosya[Denklikler]{Bolum04}{uc_ilke_denk}

\begin{dikkat}[Dikkat: ``Her $P$ için'' ifadesinin önemi]
Teorem, ilkelerin \emph{bütün önermeler} için söylenmiş hâlleri
hakkındadır. Tek bir $P$ için \lc{P ∨ ¬P}'den \lc{¬¬P → P} çıkar, ama
\lc{¬¬P → P}'den \lc{P ∨ ¬P} çıkmaz. \lc{peirce\_em} ispatında Peirce
yasasını $P$ için değil, \lc{P ∨ ¬P} önermesi için kullandığımıza dikkat
edin. Alıştırma~\ref{alis:dne-em} bu konuyu inceliyor.
\end{dikkat}

% ---------------------------------------------------------------------
\section{Bir ispat klasik mi? \texorpdfstring{\#print axioms}{\#print axioms}}\label{sec:print-axioms}

Uzun bir ispatın bir yerinde klasik mantık kullanılıp kullanılmadığını
gözle anlamak zordur. Bir teorem, klasik mantık kullanan başka bir
teoremi çağırıyor olabilir. \Lean{} bu soruyu kesin olarak cevaplar:

\leandosya[Aksiyom denetimi]{Bolum04}{print_axioms}

\lc{em\_dne} teoremi, EM'yi bir \emph{varsayım} olarak alır; hiçbir
aksiyom kullanmaz. \lc{cift\_degilleme\_eleme} ise \lc{Classical.em}'i
çağırdığı için üç aksiyoma dayanır. Bu üç aksiyom, \Lean{}'in bütün
klasik matematiğinin temelidir:

\begin{center}
\renewcommand{\arraystretch}{1.45}
\rowcolors{2}{tablosatir}{white}
\begin{tabularx}{\linewidth}{>{\raggedright\arraybackslash}p{3.3cm} X >{\centering\arraybackslash}p{1.4cm}}
\rowcolor{tabloust}
\tablobaslik{Aksiyom} & \tablobaslik{Kabaca ne söyler?} & \tablobaslik{Bölüm} \\
\lc{propext}          & Denk önermeler eşittir: \lc{(P ↔ Q) → P = Q}. & \ref{bol:aksiyomlar} \\
\lc{Quot.sound}       & Bölüm tiplerinde denk elemanlar eşittir; fonksiyon dışsallığı bundan türetilir. & \ref{bol:aksiyomlar} \\
\lc{Classical.choice} & Boş olmayan her tipten bir eleman seçilebilir (seçim aksiyomu). & \ref{bol:aksiyomlar} \\
\end{tabularx}
\end{center}

\lc{Classical.em}, \Lean{}'de bir aksiyom değil, \emph{teoremdir}: Bu üç
aksiyomdan \emph{Diaconescu teoremi} ile türetilir \citep{diaconescu1975}.
Yani \Lean{}'de klasik mantığın asıl kaynağı seçim aksiyomudur. Bu
şaşırtıcı bağlantıyı \cref{bol:aksiyomlar}'da ayrıntılı olarak
inceleyeceğiz.

% ---------------------------------------------------------------------
\section{Neden ispatlanamaz? Kripke modelleri}\label{sec:kripke-onizleme}

Şimdiye kadar $P \vee \neg P$'yi yapıcı olarak ``ispatlayamadığımızı''
söyledik. Ama bir şeyi ispatlayamamak, onun ispatlanamaz olduğunu
göstermez; belki yeterince akıllıca bir ispat bulamamışızdır. Bir
önermenin bir sistemde ispatlanamadığını göstermek için o sistemde
\emph{geçerli olup da} önermenin \emph{yanlış olduğu} bir model
bulmamız gerekir. Yapıcı mantık için bu modeller, Saul Kripke'nin
1965'te tanıttığı modellerdir \citep{kripke1965}.

\begin{tanim}{Kripke modeli (yapıcı mantık için)}{kripke}
Bir \textbf{Kripke modeli}, şunlardan oluşur:
\begin{itemize}
  \item Bir \emph{bilgi durumları} kümesi $W$ ve bu durumlar
        arasında bir sıralama $\leq$. ($w \leq v$: ``$v$, $w$'den daha
        fazla bilgi içeren bir gelecek durumdur.'')
  \item Her durumda hangi atomik önermelerin \emph{bilindiğini}
        gösteren bir ilişki. Bir kez bilinen bir atom, sonraki bütün
        durumlarda da bilinir (\emph{kalıcılık}).
\end{itemize}
Bileşik önermelerin bir $w$ durumunda bilinmesi ($w \Vdash \varphi$)
şöyle tanımlanır:
\begin{itemize}
  \item $w \Vdash P \wedge Q$ \;ancak ve ancak\; $w \Vdash P$ ve $w \Vdash Q$,
  \item $w \Vdash P \vee Q$ \;ancak ve ancak\; $w \Vdash P$ ya da $w \Vdash Q$,
  \item $w \Vdash P \to Q$ \;ancak ve ancak\; $w$'den sonraki \emph{her}
        $v$ durumunda, $v \Vdash P$ ise $v \Vdash Q$,
  \item $w \Vdash \neg P$ \;ancak ve ancak\; $w$'den sonraki \emph{hiçbir}
        durumda $P$ bilinmiyorsa.
\end{itemize}
\end{tanim}

Bu modellerin arkasındaki sezgi, zaman içinde bilgi biriktiren bir
matematikçidir. $\neg P$'yi bilmek, ``$P$'yi \emph{hiçbir zaman}
öğrenmeyeceğim'' demektir; şu an $P$'yi bilmemek bunun için yetmez.

\begin{figure}[htbp]
\centering
\begin{tikzpicture}[
    dunya/.style={circle, draw=anarenk, very thick, fill=tanimbg,
                  minimum size=1.5cm, font=\baslikfont\bfseries},
    ok/.style={-{Stealth[length=3mm]}, very thick, anarenk},
    bilgi/.style={font=\small, align=left, anchor=west},
  ]
  \node[dunya] (w0) at (0,0) {$w_0$};
  \node[dunya, fill=bittibg, draw=bittiframe] (w1) at (0,3.2) {$w_1$};
  \draw[ok] (w0) -- node[left=4pt, font=\small, text=anarenk]{$\leq$} (w1);

  \node[bilgi] at (1.2,3.2) {%
    $P$ \textcolor{bittiframe}{\textbf{bilinir}} (atom)};
  \node[bilgi] at (1.2,0.45) {%
    $P$ \textcolor{hataframe}{\textbf{bilinmez}}};
  \node[bilgi] at (1.2,-0.15) {%
    $\neg P$ \textcolor{hataframe}{\textbf{bilinmez}}: gelecekte ($w_1$) $P$ öğreniliyor};
  \node[bilgi] at (1.2,-0.75) {%
    $P \vee \neg P$ \textcolor{hataframe}{\textbf{bilinmez}}: iki taraf da bilinmiyor};
  \node[bilgi] at (1.2,-1.35) {%
    $\neg\neg P$ \textcolor{bittiframe}{\textbf{bilinir}}: hiçbir gelecek durumda $\neg P$ yok};
  \draw[black!25, rounded corners=4pt] (-1.0,-1.8) rectangle (9.6,4.15);
\end{tikzpicture}
\caption{İki durumlu bir Kripke modeli. $w_0$ durumunda $\neg\neg P$
bilinir ama $P$ bilinmez; $P \vee \neg P$ de bilinmez.}
\label{fig:kripke-em}
\end{figure}

\Cref{fig:kripke-em}'deki modelde $w_0$ durumu, ``$P$'nin doğru olup
olmadığını henüz bilmiyoruz, ama ileride öğreneceğiz'' durumudur. Bu
durumda:

\begin{itemize}
  \item $P$ bilinmez (henüz öğrenilmedi).
  \item $\neg P$ de bilinmez, çünkü gelecekteki $w_1$ durumunda $P$
        öğreniliyor. ``$P$'yi hiçbir zaman öğrenmeyeceğim'' diyemeyiz.
  \item Dolayısıyla $P \vee \neg P$ bilinmez.
  \item Buna karşılık $\neg\neg P$ bilinir: $w_0$'dan sonraki hiçbir
        durumda $\neg P$ bilinmiyor. O hâlde bu modelde
        $\neg\neg P \to P$ de bilinmez.
\end{itemize}

\begin{teorem}{Kripke modellerine göre geçerlilik}{kripke-gecerlilik}
Yapıcı önermeler mantığında ispatlanabilen her önerme, her Kripke
modelinin her durumunda bilinir \citep{kripke1965}.
\end{teorem}

Bu teoremi kabul edersek sonuç hemen çıkar: $P \vee \neg P$, bir Kripke
modelinin bir durumunda bilinmediğine göre, yapıcı olarak
ispatlanamaz. Teoremin kendisini \ciltiki{2}'de, Kripke modellerini
\Lean{}'de tanımladıktan sonra ispatlayacağız.

% ---------------------------------------------------------------------
\section{Klasik mantık yapıcı mantığın içinde: Glivenko ve Gödel}\label{sec:glivenko}

Klasik mantık yapıcı mantıktan güçlüdür. Ama şaşırtıcı biçimde, klasik
mantık yapıcı mantığın \emph{içine} gömülebilir. Bunun ilk örneğini
\cref{sec:em-yanlis-olamaz}'te gördük: $P \vee \neg P$ yapıcı
değildir, ama $\neg\neg(P \vee \neg P)$ yapıcıdır. Aynı olgu Peirce
yasası için de geçerlidir:

\leandosya[Peirce yasası, çift değilleme altında]{Bolum04}{glivenko_peirce}

\noindent Bu bir rastlantı değildir.

\begin{teorem}{Glivenko, 1929}{glivenko}
Bir $\varphi$ önermesi \emph{klasik önermeler mantığında}
ispatlanabiliyorsa, $\neg\neg\varphi$ önermesi \emph{yapıcı önermeler
mantığında} ispatlanabilir \citep{glivenko1929}.
\end{teorem}

Glivenko teoremi, klasik önermeler mantığının her totolojisinin yapıcı
mantıkta ``yanlış olamaz'' biçiminde zaten bulunduğunu söyler. Yüklemler
mantığında ise $\neg\neg$'yi yalnızca başa koymak yetmez; değillemeleri
önermenin içine yaymak gerekir. Bunu yapan dönüşüm, Gödel ve Gentzen'in
birbirinden bağımsız olarak 1933'te buldukları \emph{çift değilleme
çevirisidir} \citep{godel1933, gentzen1936}.

\begin{tanim}{Gödel--Gentzen çevirisi (önermeler için)}{godel-gentzen}
Her $\varphi$ önermesine bir $\varphi^N$ önermesi şöyle karşılık getirilir:
\begin{center}
\renewcommand{\arraystretch}{1.3}
\begin{tabular}{r@{\;\;:=\;\;}l@{\hspace{3em}}r@{\;\;:=\;\;}l}
$P^N$                 & $\neg\neg P$ \quad ($P$ atomik) &
$(\varphi \wedge \psi)^N$ & $\varphi^N \wedge \psi^N$ \\
$\bot^N$              & $\bot$ &
$(\varphi \to \psi)^N$    & $\varphi^N \to \psi^N$ \\
$(\neg\varphi)^N$     & $\neg\varphi^N$ &
$(\varphi \vee \psi)^N$   & $\neg(\neg\varphi^N \wedge \neg\psi^N)$ \\
\end{tabular}
\end{center}
O zaman $\varphi$ klasik olarak ispatlanabilir ancak ve ancak $\varphi^N$
yapıcı olarak ispatlanabilir.
\end{tanim}

\begin{figure}[htbp]
\centering
\begin{tikzpicture}[
    bolge/.style={draw=#1, fill=#1!7, very thick, rounded corners=8pt},
  ]
  \draw[bolge=anarenk] (0,0) rectangle (11,4.4);
  \node[anchor=north west, font=\baslikfont\bfseries, text=anarenk] at (0.3,4.25) {Klasik mantık};
  \draw[bolge=ikincil] (0.6,0.5) rectangle (5.6,3.3);
  \node[anchor=north west, font=\baslikfont\bfseries\small, text=ikincil] at (0.8,3.15) {Yapıcı mantık};
  \node[draw=anarenk, fill=white, rounded corners=3pt, font=\small, inner sep=5pt] (phi) at (8.6,2.0) {$\varphi$};
  \node[draw=ikincil, fill=white, rounded corners=3pt, font=\small, inner sep=5pt] (phin) at (3.1,1.9) {$\varphi^N$};
  \draw[-{Stealth[length=3mm]}, very thick, vurgu] (phi) to[bend right=20] (phin);
  \node[font=\baslikfont\scriptsize\bfseries, text=vurgu, fill=anarenk!7, inner sep=2pt] at (7.6,3.25) {çift değilleme çevirisi};
  \node[font=\baslikfont\scriptsize, text=black!60, align=center, below=4pt of phi] {klasik olarak\\ispatlanabilir};
  \node[font=\baslikfont\scriptsize, text=black!60, align=center, below=4pt of phin] {yapıcı olarak\\ispatlanabilir};
\end{tikzpicture}
\caption{Çift değilleme çevirisi, klasik mantığı yapıcı mantığın içine gömer.}
\label{fig:cevirisi}
\end{figure}

\begin{onemlifikir}[Önemli Fikir: Klasik mantık tutarlıysa, yapıcı mantık da tutarlıdır — ve tersi]
Çift değilleme çevirisi $\bot$'u $\bot$'a götürür. Dolayısıyla klasik
mantıkta bir çelişki türetilebilseydi, yapıcı mantıkta da türetilebilirdi.
Klasik mantığı kullanmak, yapıcı mantığa göre \emph{yeni bir tutarsızlık
riski} taşımaz. Brouwer ile Hilbert arasındaki tartışmanın bugünkü
barışçıl sonucunun teknik nedeni budur.
\end{onemlifikir}

% ---------------------------------------------------------------------
\section{Karar verilebilir önermeler}\label{sec:decidable}

Bazı önermeler için üçüncü hâlin olanaksızlığı yapıcı olarak doğrudur:
\emph{doğru mu yanlış mı olduğunu hesaplayabildiğimiz} önermeler için.
``$3 < 5$'' önermesinin doğru olup olmadığını bir bilgisayar programı
bize söyleyebilir. Bu tür önermelere \emph{karar verilebilir} denir.

\Lean{}'de bir önermenin karar verilebilir olması, \lc{Decidable P}
türünde bir \emph{örneğin} (İng.\ \emph{instance}) bulunmasıdır. Bu örnek,
$P$'nin doğru mu yanlış mı olduğunu hesaplayan bir algoritmadır.

\leandosya[Karar verilebilirlik]{Bolum04}{decidable}

\begin{table}[htbp]
\centering
\renewcommand{\arraystretch}{1.45}
\rowcolors{2}{tablosatir}{white}
\begin{tabular}{>{\raggedright\arraybackslash}p{3.1cm}
                >{\raggedright\arraybackslash}p{4.3cm}
                >{\raggedright\arraybackslash}p{4.3cm}}
\rowcolor{tabloust}
\tablobaslik{} & \tablobaslik{\lc{Bool}} & \tablobaslik{\lc{Prop}} \\
Elemanları         & yalnızca \lc{true} ve \lc{false}  & sonsuz sayıda önerme \\
Ne anlatır?        & bir \emph{hesaplamanın} sonucunu  & bir \emph{iddiayı} \\
``Ya ya da'' ilkesi & her zaman geçerli (iki eleman var) & yalnızca klasik mantıkla ya da \lc{Decidable} ile \\
Köprü              & \multicolumn{2}{c}{\lc{decide : (P : Prop) → [Decidable P] → Bool}} \\
\end{tabular}
\caption{\lc{Bool} ile \lc{Prop} arasındaki fark.}
\label{tab:bool-prop}
\end{table}

\leandosya[Karar verilebilir önermeler için DNE]{Bolum04}{decidable_dne}

\noindent Bu teoremin \lc{cift\_degilleme\_eleme}'den farkı, köşeli
parantez içindeki \lc{[Decidable P]} varsayımıdır. Bu varsayım
sayesinde \lc{Classical.em} yerine \lc{Decidable.em} kullanılır. Doğal
sayıların eşitliği, sıralaması, sonlu kümeler üzerindeki önermeler gibi
pek çok somut matematiksel önerme karar verilebilirdir ve onlar için
klasik akıl yürütme yapıcı olarak da geçerlidir.

% ---------------------------------------------------------------------
\section{Pratikte: Hangi mantıkla çalışmalı?}\label{sec:pratikte}

\Lean{} kullanıcısı olarak bu tartışmadan çıkarmanız gereken pratik
sonuçlar şunlardır:

\begin{itemize}
  \item \textbf{\Mathlib{} klasiktir.} \Lean{}'in matematik kütüphanesi
        klasik mantığı serbestçe kullanır. Gerçek matematik yaparken
        \lc{Classical.em}'den kaçınmanız gerekmez.
  \item \textbf{Yapıcı ispat daha fazla şey verir.} Bir ispat klasik
        mantık kullanmıyorsa, \Lean{} onu bir program olarak
        \emph{çalıştırabilir}. Bir varlık ispatından, var olan nesneyi
        hesaplayan bir program elde edebilirsiniz. Klasik ispatlar ise
        \leankw{noncomputable} olarak işaretlenmek zorunda kalabilir.
  \item \textbf{Önce yapıcı deneyin.} Bu kitabın Kısım~I ve II'sindeki
        alıştırmaların çoğu yapıcıdır. Klasik mantığa ihtiyaç
        duyduğunuzda, bunun nedenini anlamaya çalışın: Çoğu zaman
        ispatın içinde bir ``hangisi?'' sorusu saklıdır.
  \item \textbf{\lc{\#print axioms} dürüst bir tanıktır.} Bir teoremin
        hangi ilkelere dayandığından emin değilseniz, sorun.
\end{itemize}

% ---------------------------------------------------------------------
\section{Taktik ve terim özeti}

\begin{center}
\renewcommand{\arraystretch}{1.45}
\rowcolors{2}{tablosatir}{white}
\begin{tabularx}{\linewidth}{>{\raggedright\arraybackslash}p{4.6cm} X}
\rowcolor{tabloust}
\tablobaslik{Araç} & \tablobaslik{Ne yapar?} \\
\lc{Classical.em P}                 & \lc{P ∨ ¬P} (klasik). \leantac{cases} ile birlikte klasik durum ayırma. \\
\lc{Classical.byContradiction}      & \lc{(¬P → False) → P}: çelişki ile ispat (klasik). \\
\lc{Decidable.em P}                 & \lc{[Decidable P]} varken \lc{P ∨ ¬P} (yapıcı). \\
\leantac{decide}                    & Karar verilebilir bir önermeyi hesaplayarak ispatlar. \\
\lc{decide P}                       & Karar verilebilir \lc{P}'yi \lc{Bool} değerine çevirir. \\
\lc{\textcolor{leancmd}{\bfseries\#print axioms} t} & \lc{t}'nin dayandığı aksiyomları listeler. \\
\end{tabularx}
\end{center}

% ---------------------------------------------------------------------
\begin{arastirmanotu}
\textbf{Peirce yasası ve programlama.} Timothy Griffin 1990'da, Peirce
yasasının Curry--Howard eşlemesi altında programlama dillerindeki
\emph{call/cc} (``mevcut devamı yakala'') kontrol operatörünün tipine
karşılık geldiğini gösterdi \citep{griffin1990}. Böylece klasik mantık
da bir hesaplama anlamı kazandı: Klasik ispatlar, ``geri dönüp fikir
değiştirebilen'' programlardır. \Cref{sec:em-yanlis-olamaz}'teki
şeytan hikâyesi, Philip Wadler'ın bu fikri anlatmak için kullandığı
benzetmeden uyarlanmıştır \citep{wadler2003}.

\textbf{Yapıcı matematik.} Errett Bishop'un 1967 tarihli kitabı, gerçel
analizin büyük bölümünün yapıcı olarak geliştirilebileceğini gösterdi
\citep{bishop1967}. Yapıcı matematiğin felsefi ve teknik temelleri için
\citet{troelstra1988} ve \citet{dummett2000} temel kaynaklardır.

\textbf{Lean'de yapıcılık.} \Lean{}'in \lc{Prop} evreni
\emph{kanıt-ilgisizdir} (\cref{bol:bagimli}): Bir önermenin iki
kanıtı her zaman eşit sayılır. Bu tasarım, klasik aksiyomların hesaplanan
programlara ``bulaşmasını'' engeller: \lc{Prop} içinde klasik akıl
yürütme, \lc{Type} içindeki hesaplamaları bozmaz. Bu dengenin kesin
ifadesi hâlâ etkin bir araştırma konusudur \citep{carneiro2019}.
\end{arastirmanotu}

\begin{bolumozeti}
\begin{itemize}
  \item BHK yorumunda bir kanıt bir \emph{inşadır}. $P \vee Q$'nun kanıtı
        hangi tarafın doğru olduğunu söylemelidir. Bu yüzden
        $P \vee \neg P$ yapıcı olarak genel bir yasa değildir.
  \item \Lean{}'de klasik mantık \lc{Classical.em} ve
        \lc{Classical.byContradiction} ile açıkça çağrılır.
  \item EM, DNE ve Peirce yasası yapıcı mantıkta birbirine denktir.
        Herhangi birini eklemek klasik mantığı verir.
  \item \lc{\#print axioms}, bir ispatın klasik olup olmadığını gösterir.
        \lc{Classical.em}, seçim aksiyomundan türetilir.
  \item Kripke modelleri, bir önermenin yapıcı olarak ispatlanamadığını
        göstermenin aracıdır.
  \item Glivenko teoremi ve çift değilleme çevirisi, klasik mantığı
        yapıcı mantığın içine gömer.
  \item Karar verilebilir önermeler için üçüncü hâl yapıcı olarak
        geçerlidir.
\end{itemize}
\end{bolumozeti}

\begin{alistirmalar}
\alistirma{1} (Klasik) \lc{(P → Q) → ¬P ∨ Q}. Bu, implikasyonun
``maddi'' tanımıdır. (\lc{a4\_1})

\alistirma{1} (Yapıcı) Bir önceki alıştırmanın tersi:
\lc{¬P ∨ Q → (P → Q)}. Neden bu yön yapıcıdır? (\lc{a4\_2})

\alistirma{2} (Klasik) \lc{¬(P → Q) → P ∧ ¬Q}. İki bileşenden hangisi
klasik mantık gerektiriyor? \lc{\#print axioms} ile kendi
ispatınızı denetleyin. (\lc{a4\_3})

\alistirma{2} (Klasik) Kontrapozisyonun tersini terim modunda
yazın: \lc{(¬Q → ¬P) → (P → Q)}. (\lc{a4\_4})

\alistirma{2} (Klasik) \emph{Dummett ilkesi}: \lc{(P → Q) ∨ (Q → P)}.
İki rastgele önermeden biri her zaman ötekini gerektirir! Bu ilkenin
sezgisel olarak neden tuhaf olduğunu tartışın. (\lc{a4\_5})

\alistirma{2} (Yapıcı) \lc{[Decidable P]} varsayımıyla De Morgan'ın zor
yönünü ispatlayın. \lc{Decidable Q} varsayımı da gerekli midir?
(\lc{a4\_6})

\alistirma{3} (Yapıcı) Glivenko teoremini Dummett ilkesi için
doğrulayın: \lc{¬¬((P → Q) ∨ (Q → P))}. (\lc{a4\_7})

\alistirma{3}\label{alis:dne-em} (Yapıcı) Her önerme için çift
değilleme elemeyi biliyorsak üçüncü hâli de biliriz:
\lc{(∀ X : Prop, ¬¬X → X) → P ∨ ¬P}. Hangi \lc{X} için kullanmanız
gerekti? (\lc{a4\_8})

\alistirma{3} (Kâğıt-kalem) \Cref{fig:kripke-em}'deki modelde Peirce
yasasının ($Q$ yerine $\bot$ alarak) $w_0$ durumunda bilinmediğini
gösterin.

\alistirma{3} (Kâğıt-kalem) Üç durumlu bir Kripke modeli kurarak
Dummett ilkesinin yapıcı olarak ispatlanamadığını gösterin. İpucu:
$w_0$'dan iki farklı yöne dallanan $w_1$ ve $w_2$ durumları düşünün.

\alistirma{3} (Tartışma) Teorem~\ref{trm:irrasyonel-us}'in yapıcı bir ispatı
var mıdır? İpucu: $a = \sqrt{2}$, $b = 2\log_2 3$ seçimini inceleyin.
Bu seçimle ispatın hangi adımları yapıcıdır?
\end{alistirmalar}
